% =======================================================
% 2. TECNOLOGIE UTILIZZATE
% =======================================================
\section{Tecnologie Utilizzate}
Le tecnologie qui descritte confermano e consolidano le scelte motivate in dettaglio in \link{link_analisi_stato_arte}{Analisi dello Stato dell'Arte}, Sezioni 2--3, a cui si rimanda per l'analisi comparativa completa delle alternative scartate.

\subsection{Linguaggi di programmazione}
Il Frontend è sviluppato in \textbf{TypeScript}, per la tipizzazione statica che riduce gli errori di integrazione con le interfacce dei payload AI. Il Backend è sviluppato in \textbf{Python}, scelto per la maturità del suo ecosistema nell'integrazione con strumenti di intelligenza artificiale e NLP e per la conoscenza pregressa del linguaggio da parte di alcuni membri del gruppo.

\subsection{Framework}
\begin{itemize}
	\item \textbf{\link{https://vuejs.org/guide/essentials/reactivity-fundamentals.html}{Vue.js}} (Frontend): framework a componenti scelto per la progressività (si parte da codice semplice, aggiungendo complessità solo dove necessario), la separazione netta tra struttura HTML e styling CSS e la bassa curva di apprendimento, coerente con l'esigenza di un'interfaccia semplice e intuitiva anche per utenti non tecnici.
	\item \textbf{\link{https://fastapi.tiangolo.com/tutorial/dependencies/}{FastAPI}} (Backend): framework per API REST asincrono, scelto per le prestazioni elevate e il supporto nativo alle chiamate asincrone, fondamentale per non bloccare il servizio durante l'attesa delle risposte dell'LLM. Fornisce inoltre validazione dei dati integrata, generazione automatica della documentazione (Swagger/OpenAPI) e gestione di middleware per logging, sicurezza e rate limiting.
\end{itemize}

\subsection{Librerie}
\begin{itemize}
	\item \textbf{\link{https://codemirror.net/docs/guide/}{CodeMirror}}: componente editor di testo integrato nel Frontend, scelto per il bilanciamento tra leggerezza ed estensibilità. Fornisce un ecosistema di estensioni pronte per Markdown, syntax highlighting e keybinding, oltre alla possibilità di integrare comandi personalizzati per l'interazione diretta con l'LLM dalla finestra di scrittura.
	\item Librerie di validazione dati e generazione di schemi (integrate in FastAPI tramite Pydantic) per la definizione tipizzata dei payload delle operazioni AI (Sezione~\ref{sec:interfaccia-api}).
\end{itemize}

\subsection{Persistenza dei dati}
\label{sec:persistenza}
Il gruppo ha scelto di non implementare alcun sistema di persistenza server-side in questa fase: la gestione dei dati si basa esclusivamente sul file system locale del browser (salvataggio, apertura ed esportazione della nota, R75--R77), che copre i requisiti obbligatori del capitolato. La persistenza remota delle note e la gestione di utenti autenticati sono funzionalità opzionali (R4-V-Op, R78--R86) che il gruppo ha deciso di non realizzare, per concentrare le risorse disponibili sul consolidamento dei requisiti obbligatori; non è pertanto presente alcun database nello stack tecnologico né nell'architettura di deployment (Sezione~\ref{sec:arch-deployment}). Lo stato di questi requisiti opzionali è riportato in dettaglio in Sezione~\ref{sec:tracciamento}.

\subsection{Strumenti di sviluppo}
\begin{itemize}
	\item \textbf{\link{https://docs.docker.com/compose/}{Docker Compose}} per la containerizzazione e l'esecuzione coordinata di Frontend e Backend, scelto rispetto a un'installazione manuale (poco riproducibile) o a una macchina virtuale (overhead eccessivo per le esigenze del progetto); riduce le differenze tra ambienti locali e semplifica l'avvio del sistema con un unico comando.
	\item \textbf{Git} come strumento di versionamento del codice, con tracciamento delle attività tramite issue (requisito R5-Q-O).
	\item \textbf{StarUML} per la modellazione UML e la produzione dei diagrammi architetturali.
\end{itemize}

\subsection{Motivazioni delle scelte tecnologiche}
Le alternative scartate e i relativi motivi sono riportati integralmente in \link{link_analisi_stato_arte}{Analisi dello Stato dell'Arte}, Sezioni 2--3; si riepilogano qui i punti salienti:
\begin{itemize}
	\item \textbf{Angular} e \textbf{React} sono stati scartati a favore di Vue.js rispettivamente per l'eccessivo peso/complessità rispetto agli obiettivi del sistema e per il rischio di debito tecnico iniziale dato dall'assenza di standard architetturali consolidati nel team. La contropartita di questa scelta è che Vue.js, a differenza di Angular, non fornisce nativamente un sistema di Dependency Injection né impone pattern architetturali: la struttura MVC/Observer descritta in Sezione~\ref{sec:arch-frontend} è stata quindi progettata e implementata interamente dal gruppo, senza supporto diretto dal framework;
	\item \textbf{Monaco Editor} è stato scartato a favore di CodeMirror per l'elevato consumo di risorse e la configurazione più complessa;
	\item \textbf{Django} e \textbf{Flask} sono stati scartati a favore di FastAPI: Django per la natura sincrona e il peso eccessivo rispetto a un sistema che deve restare semplice, Flask per la mancanza nativa di validazione dei dati e supporto asincrono.
\end{itemize}
Un vincolo trasversale a tutte le scelte tecnologiche è l'interazione con i modelli LLM (Gemma, ChatGPT-OSS, Qwen) tramite API aderenti allo \link{https://platform.openai.com/docs/api-reference/chat/create}{standard OpenAI} (R2-V-O, R3-V-O): questo vincolo, unito alla necessità di poter sostituire il provider LLM senza impatti sul resto del sistema, ha guidato l'adozione del pattern Adapter descritto in Sezione~\ref{sec:design-pattern}.
